% Complete root/B/D actual-file audits and exact Noether source checks PASS.
\begin{proposition}[Proved: first-power Noether differents on all punctured covers]
\label{D:cover-uniform-Noether-different}
Use the complete equicharacteristic surface ring $R$, the fixed regular
Noether normalizations $S_i\subset R$, and the fixed ideal
$J\subset\sum_i\mathfrak N(R/S_i)$ from
Proposition~\ref{D:cover-uniform-canonical-defect}. For every finite normal
$R$-torsion-free algebra $A$ etale on the puncture, one has
\[
 \mathfrak N(R/S_i)A\subset\mathfrak N(A/S_i)
 \quad\hbox{for every }i,\qquad
 JA\subset\sum_i\mathfrak N(A/S_i).                     \tag{D.ND.1}
\]
Consequently, multiplication by every $z\in J$ on every maximal
Cohen--Macaulay $A$-module factors through a finite free $A$-module.
In particular $J\operatorname{Ext}^1_A(M,N)=0$ for every such source $M$
and every finite target $N$.

The same factorization holds after a finite local Artinian coefficient
extension $T/k$ for any finite $A_T=A\otimes_kT$-module $M_T$ which is
$T$-flat and whose special fibre is maximal Cohen--Macaulay. Thus this
first-power ideal also supplies the relative source-only annihilator
and prescribed-jet lifting input of
Lemma~\ref{D:relative-Noether-different-jets}.
\end{proposition}
\begin{proof}
First suppose $A$ is a domain, and fix $i$. The canonical-defect proof,
including its specified generic trace identifications, established
\[
 z\omega_{A/S_i}\subset A\tau_{A/S_i}
 \quad(z\in\mathfrak N(R/S_i)),                        \tag{D.ND.2}
\]
where $\tau$ is the trace section and is nonzero generically.
Define an $A$-linear map
\[
 \varphi_z:\omega_{A/S_i}\longrightarrow A,
 \qquad \alpha\longmapsto z\alpha/\tau_{A/S_i}.
\]
It takes values in $A$ by (D.ND.2), and
$\varphi_z(\tau_{A/S_i})=z$. Since $A$ is finite free over $S_i$,
the evaluation description of the Noether different applies and gives
$z\in\mathfrak N(A/S_i)$
\citep[Tag 0BVT, Lemma 49.6.7 and its complete proof]
{C_StacksCanonicalDefect}. This proves (D.ND.1); finite normal products
are treated componentwise.

Write an element of $\mathfrak N(A/S_i)$ as the product
$\mu(\eta)$ of a separability tensor
$\eta=\sum_\alpha a_\alpha\otimes b_\alpha\in A\otimes_{S_i}A$
annihilated by all diagonal differences. If $M$ is maximal
Cohen--Macaulay over $A$, then it is finite free over the regular ring
$S_i$. The maps
\[
 M\longrightarrow A\otimes_{S_i}M\longrightarrow M,
 \qquad m\longmapsto\sum_\alpha a_\alpha\otimes b_\alpha m,
 \qquad a\otimes m\longmapsto am
\]
are $A$-linear and have composite $\mu(\eta)\operatorname{id}_M$.
The middle term is finite free over $A$. Sums of these factorizations
prove the assertion for $J$. Lifting the second map from the finite free
middle module through an arbitrary surjection with target $M$ proves the
stated Ext annihilation.

For the relative assertion, put $S_{i,T}=S_i\otimes_kT$. The module
$M_T$ is finite over this complete local ring. It is in fact free over
$S_{i,T}$. Choose lifts of a basis of the free $S_i$-module $M_T/\mathfrak m_T M_T$;
they give a surjection $S_{i,T}^a\to M_T$ by nilpotent Nakayama. In its
kernel sequence, tensoring with $k=T/\mathfrak m_T$ remains exact on
the left because $M_T$ is $T$-flat. The kernel has zero reduction, hence
vanishes by the same Nakayama argument. Base-change the actual tensors
$\eta$ to $T$. The preceding two maps then factor multiplication through
the finite free module $A_T\otimes_{S_{i,T}}M_T$.
No flatness condition on a target module is needed for this source
factorization. Once a chosen parameter pair is regular on the target,
the two successive one-parameter divisions in the already proved
prescribed-jet lemma apply unchanged.
\end{proof}

\begin{remark}
This refines the ideal $I^3$ in
Lemma~\ref{D:cover-uniform-MCM-annihilator} in the present
equicharacteristic, perfect-residue-field setting. It does not replace
an element of a Noether different by an invertible separability
idempotent on a finite Artin quotient. In particular it supplies no
Morita equivalence and no bound with $O(nr^2)$ object coordinates.
\end{remark}

\begin{lemma}[Proved: a weighted genus estimate from geometric genus]
\label{D:exceptional-weighted-genus}
Let $Z\to\operatorname{Spec}A$ be a projective regular resolution of a
normal local surface with algebraically closed residue field. Let
$E_1,\ldots,E_v$ be distinct exceptional integral curves, put
$g_i=p_a(E_i)$ and $b_i=-E_i^2>0$, and suppose
$p_g=\dim_kH^1(Z,\mathcal O_Z)<\infty$. Then
\[
                  \sum_{i=1}^{v}\frac{g_i^2}{b_i}
                  \le 4p_g+2v.                         \tag{D.WG.1}
\]
In the punctured-cover setting of
Proposition~\ref{D:punctured-cover-linear-resolution}, this sum is
$O(n)$ with a constant depending only on the fixed base data.
\end{lemma}
\begin{proof}
For every nonzero effective exceptional cycle $D$, the vanishing of
$H^2(Z,\mathcal O_Z(-D))$ gives
$H^1(Z,\mathcal O_Z)\twoheadrightarrow H^1(D,\mathcal O_D)$.
Since $h^0(\mathcal O_D)\ge1$, it follows that $p_a(D)\le p_g$.
For $D=\sum_i t_iE_i$, cycle adjunction gives
\[
 p_a(D)=1+\sum_i\left[t_i(g_i-1)
               -\frac{b_i t_i(t_i-1)}2\right]
               +\sum_{i<j}t_it_j(E_i\cdot E_j).          \tag{D.WG.2}
\]
All terms in the last sum are nonnegative.

Set $a_i=(g_i-1)_+$. For each $g_i>1$, choose a nearest integer
$t_i$ to $a_i/b_i+1/2$; set $t_i=0$ for $g_i\le1$.
The maximum on the real line of
$q_i(t)=a_it-b_it(t-1)/2$ is
\[
 \frac{a_i^2}{2b_i}+\frac{a_i}{2}+\frac{b_i}{8}.
\]
Rounding changes it by at most $b_i/8$, so
$q_i(t_i)\ge a_i^2/(2b_i)$. If any $a_i>0$, this defines a nonzero
cycle, and (D.WG.2) implies
$\sum_i a_i^2/b_i\le2p_g$. If all $a_i=0$, that inequality is
automatic. Finally, for $g_i\ge1$,
$g_i^2\le2a_i^2+2$, and $b_i\ge1$ by integrality. The terms with
$g_i=0$ vanish, so (D.WG.1) follows. The last assertion combines the
degree-linear bound for $p_g$ in
Proposition~\ref{D:cover-uniform-canonical-defect} with the
degree-linear number of components in
Proposition~\ref{D:punctured-cover-linear-resolution}.
\end{proof}
